\chapter{2004年全国IV卷（文）}
\section{单选题}
\begin{problem}
设集合 \(U = \{0, 1, 2, 3, 4, 5\}\)，集合 \(M = \{0, 3, 5\}\)，\(N = \{1, 4, 5\}\)，则 \(M \cap (\complement_U N) =\)（\quad）

\choices
    {\(\{5\}\)}
    {\(\{0,3\}\)}
    {\(\{0, 2, 3, 5\}\)}
    {\(\{0, 1, 3, 4, 5\}\)}
\end{problem}

\begin{problem}
函数 \(y = \e^{2x}\)（\(x \in \R\)）的反函数为（\quad）

\choices
    {\(y = 2\ln x\)（\(x > 0\)）}
    {\(y = \ln(2x)\)（\(x > 0\)）}
    {\(y = \frac{1}{2}\ln x\)（\(x > 0\)）}
    {\(y = \frac{1}{2}\ln 2x\)（\(x > 0\)）}
\end{problem}

\begin{problem}
正三棱柱侧面的一条对角线长为 2，且与底面成 \(45^{\circ}\) 角，则此三棱柱的体积（\quad）

\choices
    {\(\frac{\sqrt{6}}{2}\)}
    {\(\sqrt{6}\)}
    {\(\frac{\sqrt{6}}{6}\)}
    {\(\frac{\sqrt{6}}{3}\)}
\end{problem}

\begin{problem}
函数 \( y = (x + 1)^2 (x - 1) \) 在 \( x = 1 \) 处的导数等于（\quad）

\choices
    {1}
    {2}
    {3}
    {4}
\end{problem}

\begin{problem}
为了得到函数 \( y = 3 \times \left(\frac{1}{3}\right)^x \) 的图象，可以把函数 \( y = \left(\frac{1}{3}\right)^x \) 的图象（\quad）

\choices
    {向左平移 3 个单位长度}
    {向右平移 3 个单位长度}
    {向左平移 1 个单位长度}
    {向右平移 1 个单位长度}
\end{problem}

\begin{problem}
等差数列 \(\{a_{n}\}\) 中，\(a_{1} + a_{2} + a_{3} = -24\)，\(a_{18} + a_{19} + a_{20} = 78\)，则此数列前 20 项和等于（\quad）

\choices
    {160}
    {180}
    {200}
    {220}
\end{problem}

\begin{problem}
已知函数 \( y = \log_{\frac{1}{4}} x \) 与 \( y = kx \) 的图象有公共点 \(A\). 且点 \(A\) 的横坐标为 2，则 \( k = \)（\quad）

\choices
    {\(-\frac{1}{4}\)}
    {\(\frac{1}{4}\)}
    {\(-\frac{1}{2}\)}
    {\(\frac{1}{2}\)}
\end{problem}

\begin{problem}
已知圆 \(C\) 的半径为 2，圆心在 \(x\) 轴的正半轴上，直线 \(3x + 4y + 4 = 0\) 与圆 \(C\) 相切，则圆 \(C\) 的方程为（\quad）

\choices
    {\(x^{2} + y^{2} - 2x - 3 = 0\)}
    {\(x^{2} + y^{2} + 4x = 0\)}
    {\(x^{2} + y^{2} + 2x - 3 = 0\)}
    {\(x^{2} + y^{2} - 4x = 0\)}
\end{problem}

\begin{problem}
从 5 位男教师和 4 位女教师中选出 3 位教师，派到 3 个班担任班主任（每班 1 位班主任）.要求这 3 位班主任中男，女教师都要有，则不同的选派方案共有（\quad）

\choices
    {210 种}
    {420 种}
    {630 种}
    {840 种}
\end{problem}

\begin{problem}
函数 \(y = 2\sin \left(\frac{\pi}{3} - x\right) - \cos \left(\frac{\pi}{6} + x\right)\)（\(x \in \R\)）的最小值等于（\quad）

\choices
    {\(-3\)}
    {\(-2\)}
    {\(-1\)}
    {\(-\sqrt{5}\)}
\end{problem}

\begin{problem}
已知球的表面积为 \(20\pi\)，球面上有 \(A\)，\(B\)，\(C\) 三点. 如果 \(AB = AC = 2\)，\(BC = 2\sqrt{3}\)，则球心到平面 \(ABC\) 的距离为（\quad）

\choices
    {1}
    {\(\sqrt{2}\)}
    {\(\sqrt{3}\)}
    {2}
\end{problem}

\begin{problem}
\(\triangle ABC\) 中，\(a\)，\(b\)，\(c\) 分别为 \(\angle A\)，\(\angle B\)，\(\angle C\) 的对边. 如果 \(a\)，\(b\)，\(c\) 成等差数列，\(\angle B = 30^{\circ}\)，\(\triangle ABC\) 的面积为 \(\frac{3}{2}\)，那么 \(b =\)（\quad）

\choices
    {\(\frac{1 + \sqrt{3}}{2}\)}
    {\(1 + \sqrt{3}\)}
    {\(\frac{2 + \sqrt{3}}{2}\)}
    {\(2 + \sqrt{3}\)}
\end{problem}

\section{填空题}
\begin{problem}
\(\left(x - \frac{1}{\sqrt{2}}\right)^8\) 展开式中 \(x^5\) 的系数为\fillinblank{}.
\end{problem}

\begin{problem}
已知函数 \(y = \frac{1}{2} \sin \frac{x + \pi}{A}\)（\(A > 0\)）的最小正周期为 \(3\pi\)，则 \(A =\)\fillinblank{}.
\end{problem}

\begin{problem}
向量 \(\bs{a}\)，\(\bs{b}\) 满足 \(\left(\bs{a} - \bs{b}\right) \cdot \left(2\bs{a} + \bs{b}\right) = -4\)，且 \(|\bs{a}| = 2\)，\(\left|\bs{b}\right| = 4\)，则 \(\bs{a}\) 与 \(\bs{b}\) 夹角的余弦值等于\fillinblank{}.
\end{problem}

\begin{problem}
设 \(x\)，\(y\) 满足约束条件：\(\left\{ \begin{array}{l} x + y \leqslant 1, \\ y \leqslant x, \\ y \geqslant 0, \end{array} \right.\) 则 \(z = 2x + y\) 的最大值是\fillinblank{}.
\end{problem}

\section{解答题}
\begin{problem}
已知 \(\alpha\) 为第二象限角，且 \(\sin \alpha = \frac{\sqrt{15}}{4}\)，求 \(\frac{\sin (\alpha + \frac{\pi}{4})}{\sin 2\alpha + \cos 2\alpha + 1}\) 的值.
\end{problem}

\begin{problem}
已知数列 \(\{a_{n}\}\) 为等比数列，\(a_{2} = 6\)，\(a_{5} = 162\).

\begin{enumerate}
\item 求数列 \(\{a_{n}\}\) 的通项公式；
\item 设 \(S_{n}\) 是数列 \(\{a_{n}\}\) 的前 \(n\) 项，证明 \(\frac{S_{n} \cdot S_{n+2}}{S_{n+1}^{2}} \leqslant 1\).
\end{enumerate}
\end{problem}

\begin{problem}
已知直线 \(l_{1}\) 为曲线 \(y = x^{2} + x - 2\) 在点 \((1,0)\) 处的切线，\(l_{2}\) 为该曲线的另一条切线，且 \(l_{1} \perp l_{2}\).

\begin{enumerate}
\item 求直线 \(l_{2}\) 的方程；
\item 求由直线 \(l_{1}\)，\(l_{2}\) 和 \(x\) 轴所围成的三角形的面积.
\end{enumerate}
\end{problem}

\begin{problem}
某同学参加科普知识竞赛，需回答 3 个问题. 竞赛规则规定：答对第一，二，三问题分别得 100 分，100 分，200 分，答错得零分. 假设这名同学答对第一，二，三个问题的概率分别为 0.8，0.7，0.6，且各题答对与否相互之间没有影响.

\begin{enumerate}
\item 求这名同学得 300 分的概率；
\item 求这名同学至少得 300 分的概率.
\end{enumerate}
\end{problem}

\begin{problem}
如图，四棱锥 \(P - ABCD\) 中，底面 \(ABCD\) 为矩形，\(AB = 8\)，\(AD = 4\sqrt{3}\)，侧面 \(PAD\) 为等边三角形，并且与底面所成二面角为 \(60^{\circ}\).

\begin{enumerate}
\item 求四棱锥 \(P - ABCD\) 的体积；
\item 证明：\(PA \perp BD\).
\end{enumerate}

\begin{center}
\bitmapfigure[max width=0.42\linewidth]{img/2004/national_paper_4_liberal/q21_fig1.png}
\end{center}
\end{problem}

\begin{problem}
双曲线 \(\frac{x^2}{a^2} - \frac{y^2}{b^2} = 1\)（\(a > 1\)，\(b > 0\)）的焦点距为 \(2c\)，直线 \(l\) 过点 \((a,0)\) 和 \((0,b)\)，且点 \((1,0)\) 到直线 \(l\) 的距离与点 \((-1,0)\) 到直线 \(l\) 的距离之和 \(s \geqslant \frac{4}{5} c\). 求双曲线的离心率 \(e\) 的取值范围.
\end{problem}
